% NAVCORE-SoOP — order to first flight
% Compile: tectonic navcore-runbook.tex   (or pdflatex, twice, for the ToC)
\documentclass[11pt,a4paper]{article}

\usepackage[margin=20mm,top=18mm,bottom=20mm]{geometry}
\usepackage[T1]{fontenc}
\usepackage{booktabs}
\usepackage{longtable}
\usepackage{array}
\usepackage{tabularx}
\usepackage{enumitem}
\usepackage{ragged2e}
\usepackage{xcolor}
\usepackage{titlesec}
\usepackage{fancyhdr}
\usepackage{hyperref}
\usepackage{microtype}

\definecolor{navy}{RGB}{20,50,95}
\definecolor{warnred}{RGB}{165,30,30}
\definecolor{okgreen}{RGB}{20,105,60}
\definecolor{sidegrey}{RGB}{95,100,110}

\hypersetup{
  colorlinks=true,
  linkcolor=navy,
  urlcolor=navy,
  pdftitle={NAVCORE-SoOP: order to first flight},
  pdfsubject={Ordering, assembly and bring-up runbook}
}

\titleformat{\section}{\normalfont\Large\bfseries\color{navy}}{\thesection}{0.7em}{}
\titleformat{\subsection}{\normalfont\large\bfseries\color{navy}}{\thesubsection}{0.6em}{}
\titleformat{\subsubsection}{\normalfont\normalsize\bfseries}{\thesubsubsection}{0.5em}{}

\pagestyle{fancy}
\fancyhf{}
\renewcommand{\headrulewidth}{0.3pt}
\fancyhead[L]{\footnotesize\color{sidegrey}NAVCORE-SoOP --- order to first flight}
\fancyhead[R]{\footnotesize\color{sidegrey}\thepage}

\newcommand{\file}[1]{\texttt{\small #1}}
\newcommand{\cmd}[1]{\texttt{\small #1}}
\newcommand{\flag}[1]{\textcolor{warnred}{\textbf{#1}}}
\newcommand{\good}[1]{\textcolor{okgreen}{\textbf{#1}}}
\newcommand{\degC}{$^\circ$C}

% Table columns that wrap rather than stretch, so long \file{} tokens do not
% produce justified rivers or overfull boxes.
\newcolumntype{P}[1]{>{\RaggedRight\arraybackslash}p{#1}}
\renewcommand{\arraystretch}{1.14}

\setlist[itemize]{leftmargin=1.3em,itemsep=1pt,topsep=2pt}
\setlist[enumerate]{leftmargin=1.6em,itemsep=1pt,topsep=2pt}

\begin{document}

\begin{titlepage}
\centering
\vspace*{25mm}
{\Huge\bfseries\color{navy} NAVCORE-SoOP\par}
\vspace{4mm}
{\Large Order to first flight\par}
\vspace{2mm}
{\large The complete runbook, with an audit of what is still manual,\\ soldered by hand, or not yet measured\par}
\vspace{12mm}
\rule{0.72\textwidth}{0.4pt}
\vspace{6mm}

\begin{tabular}{@{}ll@{}}
\toprule
Board & 45.1 $\times$ 46.1 mm, 6-layer, STM32H743VIT6 \\
Gate verdict & \good{READY TO ORDER} --- all 43 gated prerequisites proven \\
Hardware verdict & \flag{NOT READY TO FLY} --- 9 items need hardware \\
Toolchain & KiCad 9.0.8, ArduPilot Copter-4.7.0, Python 3 + pcbnew \\
\bottomrule
\end{tabular}

\vspace{8mm}
\begin{minipage}{0.82\textwidth}
\small\itshape
Every figure in this document was re-derived from the repository's own tooling on
19 September 2026 rather than copied from another document. Where a number is a
judgement, an assumption or genuinely unknown, it says so. Sections marked
\flag{flag} are the ones capable of costing money, a board, or a bench session.
\end{minipage}

\vfill
{\small Generated from the NAVCORE-SoOP working tree\par}
\end{titlepage}

\tableofcontents
\newpage

% =============================================================================
\section{How to use this, and how to verify it}

This is one document in the order you actually do things: order the board, order the
aircraft parts, assemble, power it up, and fly it. It exists because the same information
was spread across \file{docs/ORDER.md}, \file{docs/ASSEMBLY.md}, \file{docs/BUILD.md} and
\file{docs/BUYING.md}, and those four had drifted apart --- in three cases far enough apart
to send you shopping for parts with nowhere to put them.

\textbf{Verify before you spend. Read the verdict, not a summary:}

\begin{quote}
\cmd{python3 tools/preflight.py} \quad$\rightarrow$\quad must print \good{READY TO ORDER}
\end{quote}

The verdict is \emph{derived} from a manifest of 43 gated prerequisites, so an unrun tool,
a non-zero exit, a failed check or an unclassified manifest entry is a blocking failure.
It cannot be typed by hand and it cannot shrink silently. What it proves is that the
\emph{design} is orderable. It cannot prove anything about money, stock, or JLCPCB's own
upload preview --- those are Sections~\ref{sec:upload} onward.

\subsection{The four commands that matter}

\begin{center}
\begin{tabular}{@{}ll@{}}
\toprule
\cmd{python3 tools/preflight.py} & the gate; must say READY TO ORDER \\
\cmd{bash tools/make\_order\_bundle.sh} & regenerate \file{fab/} and zip the order \\
\cmd{bash tools/finish.sh} & the whole design pipeline, ending in the fab outputs \\
\cmd{bash tools/build\_firmware.sh} & build \file{arducopter.apj}, pinned to Copter-4.7.0 \\
\bottomrule
\end{tabular}
\end{center}

\subsection{Freshness}

\file{fab/stock-snapshot.json} is re-fetched from \file{jlcpcb.com} and
\file{tools/check\_stock.py} blocks the gate once it is older than \textbf{7 days}. It was
last refreshed 2026-09-19 and resolves \textbf{56 of 56} non-DNP codes with cover for two
boards. Everything else in \file{fab/} is regenerated by
\file{tools/make\_order\_bundle.sh}, which refuses to build the bundle if the copper layer
count is not 6.

% =============================================================================
\section{Part 1 --- Order the bare board and its assembly}
\label{sec:order-board}

\subsection{Regenerate, do not trust}
\label{sec:regenerate}

\begin{quote}
\cmd{bash tools/make\_order\_bundle.sh}
\end{quote}

This re-exports the gerbers and drill, re-runs \file{gen\_bom.py} for \emph{all three}
variant pairs, checks the layer count, and writes
\file{fab/NAVCORE-SoOP-order-bundle.zip}.

\flag{Do not hand-copy or hand-edit anything in \file{fab/}.} A \cmd{cp} of a stale variant
into the bundle is exactly how the Economic BOM once shipped quoting \file{Y2}'s
superseded, unbuyable part code for three days --- and \file{fab/ORDER.md} was recommending
those files at the checkout while it did.

\subsection{What to upload}

\begin{center}
\begin{tabular}{@{}ll@{}}
\toprule
\textbf{File} & \textbf{Where} \\
\midrule
\file{fab/gerbers/} --- zip the \emph{contents} & PCB gerbers + Excellon drill \\
\file{fab/BOM-NAVCORE-SoOP.csv} & assembly BOM, the default build \\
\file{fab/CPL-NAVCORE-SoOP.csv} & pick-and-place, the default build \\
\bottomrule
\end{tabular}
\end{center}

Or upload the bundle zip, which contains those plus \file{HOW-TO-ORDER.txt}.
\file{fab/gerbers/} holds \textbf{16 files}: 14 \file{.gbr} (6 copper, 2 paste, 2
silkscreen, 2 mask, edge cuts, drill map), \file{NAVCORE-SoOP.drl}, and
\file{NAVCORE-SoOP-job.gbrjob}.

\begin{quote}
\flag{Stop-condition at upload.} If the preview shows \textbf{four} copper layers, abort.
The power planes are missing and the board will not work. It must show six.
\end{quote}

\subsection{Fabrication settings}

\begin{center}
\begin{tabular}{@{}lll@{}}
\toprule
\textbf{Setting} & \textbf{Value} & \textbf{Why} \\
\midrule
Layers & \textbf{6} & In1 is a solid GND plane; In4 carries the rails \\
Dimensions & 45.1 $\times$ 46.1 mm & \\
Thickness & 1.6 mm & 30.5 mm stack standard \\
Outer copper & \textbf{1 oz} & every current figure assumes it \\
Surface finish & ENIG & 0402s and LGA-14s want a flat finish \\
Min track / spacing & \textbf{0.1016 mm (4 mil)} & should be free tier \\
Min via / drill & \textbf{0.45 mm / 0.20 mm} & \\
Impedance control & No & USB is a plain differential pair \\
Stencil & not needed & everything machine-placed but Section~\ref{sec:solder} \\
\bottomrule
\end{tabular}
\end{center}

The 1 oz line is not cosmetic: \flag{every current-capacity figure in
\file{check\_build.py} assumes it}, and 0.5 oz halves each trace's rating.

\subsection{Quantity, and the tier decision}

\textbf{Order 5 bare boards and 2 assembled.} Five is the multilayer minimum; two is the
usual SMT assembly minimum.

\begin{center}
\begin{tabularx}{\textwidth}{@{}lXX@{}}
\toprule
 & \textbf{Standard (panelised)} & \textbf{Economic} \\
\midrule
Board minimum & \textbf{70 $\times$ 70 mm}, so 2 $\times$ 2 panel + rails & 10 $\times$ 10 mm, no rails \\
Places & everything, including \file{U3} & everything except \file{U3} \\
Feeder fees & included & \textbf{GBP 40--46} still charged \\
Hand-fitting & none & \file{U3} (ICM-42605, second IMU) \\
\bottomrule
\end{tabularx}
\end{center}

\textbf{Recommendation: Standard, panelised 2 $\times$ 2, ``panel by JLCPCB''.} Do not
commit a panelised board file --- the board files stay exactly as gated, so nothing has to
go through DRC, connectivity and the whole gate again. A hand-panelised
\file{.kicad\_pcb} would have to, for no benefit.

If you take Economic instead, \flag{take its BOM and CPL as a pair}
(\file{*-economic.csv} for both). Mixing an Economic BOM with the default CPL hands JLC a
placement list for parts the BOM does not order. \file{tools/check\_variants.py} --- now a
gated check --- fails if any pair disagrees with the board.

\subsection{The variant files, and why two of them are identical}

\begin{center}
\begin{tabular}{@{}lrrl@{}}
\toprule
\textbf{Variant} & \textbf{BOM lines} & \textbf{Placements} & \textbf{Difference from default} \\
\midrule
\file{NAVCORE-SoOP} (default) & 68 & \textbf{129} & --- \\
\file{NAVCORE-SoOP-economic} & 68 & \textbf{128} & omits \file{U3} only \\
\file{NAVCORE-SoOP-nofpv} & 68 & \textbf{129} & \textbf{none --- byte-identical} \\
\bottomrule
\end{tabular}
\end{center}

\file{-nofpv} is byte-identical to the default pair because the 9 V VTX block is DNP
\emph{either way} (\file{design.POPULATE\_VTX} is \cmd{False}). It is a label, not a
saving. Choosing Economic only makes sense if you want to hand-fit the second IMU.

\subsection{Parts worth checking before you pay}

Stock is no longer quoted in prose. It lives in the dated snapshot
\file{fab/stock-snapshot.json}, age-gated at 7 days by \file{tools/check\_stock.py}.

Two substitutions are \textbf{not optional}:

\begin{itemize}
\item \flag{\file{Y1} must be \file{C2682774}} --- a \emph{passive crystal}, CL 20 pF. The
visually identical \file{C2901556} is an \textbf{active oscillator}: same SMD3225-4P land
pattern, but pin 4 is VDD, which this board grounds. Substituting it gives a board that
never starts.
\item \flag{\file{Y2} must be \file{C5563878}.} Its predecessor \file{C22381771} was down
to \textbf{1 unit} with no published restock date when \file{check\_stock.py} was written.
\end{itemize}

\subsection{The DNP list --- derived, never typed}

\textbf{14 designators ship unpopulated:}
\file{C68 C69 C70 C73 L5 Q3 R15 R40 R41 R42 R43 R45 U11 U18}.

\begin{center}
\begin{tabular}{@{}p{0.30\textwidth}p{0.62\textwidth}@{}}
\toprule
\textbf{Block} & \textbf{Why it is not fitted} \\
\midrule
The 9 V VTX buck (\file{L5}, \file{Q3}, \file{U18}, \file{C68}--\file{C73},
\file{R41}--\file{R45}) &
\file{POPULATE\_VTX} is \cmd{False} because its switch node could not be closed on this
placement. A VTX runs from +5 V instead. \\
\file{U11}, the CAN transceiver &
a \emph{thermal} decision, not a CAN one: it was fitted but never enabled in
\file{defaults.parm}. Removing it dropped \file{U9}'s peak from 445 mA / 179\degC{} to
375 mA / 157\degC{}. \\
\bottomrule
\end{tabular}
\end{center}

Both have pads, and the gate prints \file{[NOT FITTED]} for each on every run. Do not type
this list --- \file{make\_order\_bundle.sh} derives it from the BOM into
\file{HOW-TO-ORDER.txt}.

\subsection{Check these on screen at upload}
\label{sec:upload}

\begin{itemize}
\item Outline and board size: \textbf{45.1 $\times$ 46.1 mm}
\item Copper layer count: \textbf{6}
\item Drill count and the 0.20 mm minimum
\item \textbf{CPL orientation against their preview render} --- the one thing no local tool
substitutes for
\item \textbf{1 oz outer copper}, explicitly
\item \textbf{Panel by JLCPCB.} Standard assembly requires a single PCB of at least
$70 \times 70$~mm and this board is $45.1 \times 46.1$~mm, so the panel and its breakaway
rails have to come from somewhere. Their own panelisation is the low-risk answer: you then
confirm the panel outline in the same preview as the CPL.
\end{itemize}

\subsection{Paying}

\begin{itemize}
\item \textbf{Coupons are once-per-order} and cannot be combined. Take the SMT one.
\item \textbf{6-layer promotions recur}; worth waiting a few days if you are not in a hurry.
\item \textbf{UK import VAT is 20\%}, collected by the courier above GBP 135 with an
£8--12 handling fee. It is not a JLCPCB charge and is easy to forget when comparing
against a UK-stocked alternative.
\end{itemize}

% =============================================================================
\section{Part 2 --- Order the aircraft parts}

Order these at the same time; lead times run 2--4 weeks.

\begin{longtable}{@{}p{0.42\textwidth}p{0.10\textwidth}p{0.42\textwidth}@{}}
\toprule
\textbf{Item} & \textbf{Qty} & \textbf{Note} \\
\midrule
\endfirsthead
\toprule
\textbf{Item} & \textbf{Qty} & \textbf{Note} \\
\midrule
\endhead
\bottomrule
\endfoot
TBS Source One V5 7in DC frame & 1 & 320 mm wheelbase, 6.0 mm arms; open-source DXF \\
SpeedyBee BLS 60A 4-in-1 ESC & 1 & 30.5 $\times$ 30.5 stack \\
2806.5 1300KV motors & 4 & \textbf{19 $\times$ 19} bolt pattern, M5 shaft \\
Gemfan 7040 props & 4 + spares & 5.0 mm bore \\
Zeee 4S 6500 mAh pack & 1+ & 138 $\times$ 47 $\times$ 48 mm, 615 g, \textbf{EC5} \\
EC5-to-XT60 adapter & 2 & the pack is EC5, everything else is XT60 \\
M10 GNSS + QMC5883L compass & 1 & comes with a JST-GH 6-pin lead \\
ELRS receiver, \textbf{ESP-based} & 1 & an STM32 ELRS receiver cannot carry MAVLink \\
M3 $\times$ 14 socket screws & 16 & motors, skid sandwiched \\
M3 $\times$ 30 socket screws & 4 & \textbf{the stack bolts} \\
M3 female-female standoffs, \textbf{12 mm} & 4 & \textbf{the ESC-to-FC spacer} \\
\textbf{M3 socket assortment} (10/12/14/16/20) + washers & 1 kit & \textbf{the anti-reorder kit} \\
Frame kit's own \textbf{30 mm / 22 mm} standoffs & --- & \textbf{already in the kit; do not buy more} \\
TPU skids, printed & 4 & 3.5 mm thick, 40 mm drop, \textbf{19 $\times$ 19} holes \\
JST-SH 1.0 mm 4-pin pigtails & pack of 10 & \file{J5}, then \file{J9}/\file{J11} \\
JST-GH 1.25 mm 6-pin spare lead & 1 & closes the \file{J3} plug risk \\
5 V \textbf{passive} piezo buzzer & 1 & passive, not active \\
microSD card & 1 & 32 GB \\
Bench PSU with current limit & 1 & Section~\ref{sec:T1} does not work without one \\
Multimeter with continuity & 1 & for buzzing the ESC cable \\
Smoke stopper (XT60) & 1 & first battery connection only \\
Thermocouple or IR thermometer & 0--1 & \textbf{only} for \file{U8} now; \file{U19} logs \file{U9} \\
ST-Link V2 + fine wires or pogo pins & 1 & for the SWD bootloader \\
Nooelec SAWbird+ IR (1620 MHz) & 1 & own this by T3b, not later \\
Nooelec 1620 MHz patch antenna & 1 & own this by T3b \\
RTL-SDR (R820T2) & 1 & the T3b reference receiver \\
\end{longtable}

\flag{Do not buy a PMW3901 or a VL53L1X.} Both were deleted from the design in the Rev B
re-layout and have no pads. \file{docs/BUYING.md} told you to hand-fit them until
2026-09-19, and priced a ``GBP 227 hand-fit route'' around two parts that cannot be bought
a place for.

\subsection{The one order --- nothing here should make you order twice}

Every ``measure this, then use a different part'' branch in this document is closed by
buying the alternatives up front. The premium is about \textbf{GBP 6}.

\begin{center}
\begin{tabular}{@{}P{0.32\textwidth}P{0.56\textwidth}@{}}
\toprule
\textbf{What to add} & \textbf{What it closes} \\
\midrule
\textbf{M3 socket assortment} + washers &
Both derived joints round \emph{up} past their engagement target (4.5 mm and 6.7 mm into
the thread) and the skid's tapped depth is the only measurement that can invalidate the
length you picked. \\
\textbf{4 $\times$ 12 mm female-female spacers} &
the ESC-to-FC separation is 12.1 mm and \file{fasteners.py} says plainly that a grommet
cannot hold it. Without these the stack cannot be assembled at all. \\
the kit's own standoffs &
do \emph{not} buy 35 mm standoffs for the top plate; it has no 30.5 mm holes. \\
spare \textbf{JST-GH 6-pin} lead &
the residual \file{J3} risk is a \emph{bulky moulded boot} on the lead you receive, not the
board: the 5.30 mm clear run meets its derived 3.6 mm requirement. Test-fit the plug before
soldering; a second lead is then the fallback rather than an interruption. \\
\textbf{10} JST-SH pigtails, not 4 &
the first one you crimp badly is then not a project stall. \\
\textbf{SAWbird+ and RTL-SDR} with Tier 1 &
T3b cannot be started without them, and the board does not supply them. \\
\bottomrule
\end{tabular}
\end{center}

Three things genuinely need the parts in hand, and \textbf{none of them is an order}: the
frame kit's top-plate standoff spacing (the kit supplies it), the camera mount (deferred),
and the Radxa Zero 3W's hole pattern (deferred, and it is a \textbf{3D print}, not a
purchase). \file{tools/fasteners.py} prints this list itself under
\emph{``not derived --- these need the parts in hand''}.

% =============================================================================
\section{Part 3 --- While you wait (free, and the highest-value work)}

\begin{enumerate}
\item \textbf{Register ArduPilot board ID 7180.} Not needed to order; needed before
sharing. The request is written out in \file{firmware/BOARD-ID-REQUEST.md}, with the
registry check, the alternates and the issue text. \flag{Do not request 9001} --- it is
free but it is above the \#7199 ceiling that \file{board\_types.txt} explicitly asks you
not to cross, while 4544 free IDs sit below it.
\item \textbf{Read \file{docs/ASSEMBLY.md} Stages 3 and 4 twice.} The two screw-depth
warnings are the ones that strip a frame.
\item \textbf{Print \file{docs/img/padmap.svg} and the F.Silkscreen gerber 1:1 on paper.}
Nine pads arrive with no label (Section~\ref{sec:T3d}).
\item \textbf{Decide the USB edge orientation before cutting any loom.} All four edges
carry a connector; there may be exactly one orientation that works.
\item \textbf{Dry-run the SITL suite} if you have a Linux box:
\cmd{bash sitl/run\_scenarios.sh}.
\end{enumerate}

% =============================================================================
\section{Part 4 --- On arrival: inspect before you spend more}

\subsection{Visual inspection}

Under magnification, look at \file{U1} (\textbf{LQFP-100, 0.50 mm pitch}),
\file{U2}/\file{U3} (\textbf{LGA-14}, no accessible leads), and \file{U13}
(\textbf{TQFN-28 with an exposed pad}). You are looking for bridges and tombstones.
JLC is good at this. Look anyway.

\subsection{Confirm the machine-readable facts}

\begin{quote}
\cmd{python3 tools/preflight.py} \quad never mind the board --- confirm DRC, ERC, gerbers
and BOM/CPL are all current\\
\cmd{python3 tools/check\_variants.py} \quad every variant pair matches the board
\end{quote}

\subsection{The frame measurements that settle a real decision}

\begin{center}
\begin{tabular}{@{}p{0.40\textwidth}p{0.52\textwidth}@{}}
\toprule
\textbf{Check with calipers} & \textbf{Parsed value} \\
\midrule
top plate width & \textbf{42.50 mm} --- the battery is 47 mm and already overhangs 2.25 mm
each side; there is no room for an upward ToF beside it \\
bottom plate M3 positions & x = $\pm$14.60 and $\pm$11.00, \textbf{no 30.5 mm pattern} --- so
the FC cannot bolt under the bottom plate; that idea is closed \\
which plate is top & \textbf{inferred} from length (160.26 mm). The one genuinely open row \\
\bottomrule
\end{tabular}
\end{center}

\subsection{The three things no drawing can answer}

\begin{itemize}
\item \textbf{microSD withdrawal clearance} --- depends how \emph{you} mount it
\item \textbf{grommet compression under load} --- a property of the rubber
\item \textbf{the real TPU gear's shape, hole pattern and strength} --- the assembly is
modelled as four legs at the motor pattern on assumed dimensions
\end{itemize}

% =============================================================================
\section{Part 5 --- Mechanical assembly}

\subsection{Stack height}

\textbf{22.3 mm against 25 mm of inner space --- 2.7 mm spare.} Two figures exist and both
are correct about different things: \file{check\_build.py} gives \textbf{19.8 mm} for bare
parts, and \file{check\_mechanical.py} gives \textbf{22.3 mm} including \file{J3}'s mated
plug. \textbf{22.3 mm is the one the top plate must clear.}

\subsection{The two screw warnings --- read both}

\begin{quote}
\cmd{python3 tools/fasteners.py}
\end{quote}

\begin{itemize}
\item \flag{Motor to arm, skid sandwiched.} Needs 13.5 mm, so order M3 $\times$ 14. But
rounding 13.5 $\rightarrow$ 14 drives \textbf{5.5 mm into a thread designed for 4.0 mm}. A
motor boss is often only about 4 mm deep, and \textbf{a bottomed-out screw feels tight
while clamping nothing} --- that is the failure that throws a motor in flight. Measure the
tapped depth before fitting sixteen screws. M3 $\times$ 12 plus a washer is the usual fix.
\item \flag{FC to ESC to frame, the 30.5 mm stack bolt.} Needs \textbf{27.3 mm}, so M3
$\times$ 30. Rounding 27.3 $\rightarrow$ 30 puts \textbf{6.7 mm} into the press nut. If the
tapped depth is shallower than that the screw bottoms out and clamps nothing.
\end{itemize}

\textbf{This is the whole reason the buy list carries an M3 assortment rather than one
length.} \file{free\_island} is not involved; the point is that both joints are decided by
a measurement you have not taken yet, so own every plausible length.

\subsection{Screw heads against copper --- 0.08 mm}

\textbf{No copper is under any screw head, and the margin is thin.} The closest is
\file{TP20} at \textbf{2.83 mm} against a \textbf{2.75 mm} cap-head radius ---
\flag{0.08 mm}, and that assumes \textbf{no washers}. A washer, or a head even 0.1 mm over
nominal diameter, puts metal on copper. If you add washers, re-measure before trusting
this line.

\subsection{Sequence}

\begin{enumerate}
\item \textbf{Inspect the boards} (Section 4.1).
\item \textbf{Print the fab drawing and silkscreen 1:1.} Nine unlabelled pads ---
Section~\ref{sec:T3d}.
\item \textbf{Dry-fit the FC into the centre plate.} Confirm nothing blocks the bottom
edge where the microSD ejects.
\item \textbf{Buzz the ESC cable} both ends against the \file{J2} pinout below, before
anything is powered.
\item \textbf{Sit a skid under each arm}, holes aligned. Motor on top, wires inboard.
Thread all four screws by hand before tightening any; tighten diagonally.
\item \textbf{The stack, bottom to top: ESC, grommets, FC.} Fit the silicone grommets into
the FC's four mounting holes; ESC into the frame on its 30.5 mm pattern; FC on top with
USB-C facing the direction you will actually plug from.
\item \textbf{4 $\times$ M3 $\times$ 30 through the stack}, with the four 12 mm
female-female spacers setting the ESC-to-FC separation.
\item \textbf{30 mm standoffs for the top plate} --- from the frame kit.
\end{enumerate}

\subsection{\file{J2}, the ESC connector --- confirmed rather than assumed}

Wired \textbf{1 GND, 2 VBAT, 3 M1, 4 M2, 5 M3, 6 M4, 7 CUR, 8 TEL} (pins 9/10 are the
shell tabs). Betaflight's board page for the SpeedyBee F405 V4 gives the same FC-side
order, so it \textbf{matches} --- no cable rework needed.

The vendor table reads oddly because it lists the two \emph{ends} of the cable in opposite
order, since they face each other. Reverse the ESC column and every pair lines up:
\file{GND}$\leftrightarrow$\file{GND}, \file{BAT}$\leftrightarrow$\file{VBAT},
\file{M1}$\leftrightarrow$\file{S1}, \ldots, \file{CUR}$\leftrightarrow$\file{CURRENT}.

\textbf{Buzz it anyway.} One automated read of a wiki table is not the same as measuring
the cable you received. An 8-pin JST-SH is re-pinnable with a fine pick --- a mismatch
costs a cable rework, not a respin.

\begin{quote}
\textbf{Pin 8 is inert.} The ESC end is \file{N/A} --- stock BLHeli\_S has no telemetry
output at all. \file{SERIAL5\_PROTOCOL 16} and the \file{UART8\_RX} wire are correct and
cost nothing, but they will receive nothing until a telemetry-capable ESC (BLHeli\_32,
AM32) is fitted. That is why the board ships \file{INS\_HNTCH\_MODE 1} (throttle) rather
than 3 (RPM).
\end{quote}

% =============================================================================
\section{Part 6 --- Soldering: four joints minimum to fly}
\label{sec:solder}

Everything else plugs in.

\begin{itemize}
\item \textbf{Buzzer} --- \file{PZ1} (+) and \file{PZ2} ($-$). Must be a \textbf{passive}
piezo: PA15 is a timer channel and ArduPilot generates the tone patterns itself. An active
buzzer has its own oscillator and loses every arming and failsafe pattern.
\item \textbf{LED strip, optional} --- \file{PL1} = data, \file{PL2} = +5 V,
\file{PL3} = GND. \flag{\file{PL3} has no label.} 5 V WS2812B.
\item \flag{Do not solder to \file{P41}.} It is \textbf{5 V} and unlabelled. Identify it on
the printed gerber first, never by counting from a neighbour.
\end{itemize}

% =============================================================================
\section{Part 7 --- T1: first power, before the MCU sees anything}
\label{sec:T1}

\flag{USB disconnected. No battery. Do not skip this.} A board that goes straight to USB
can destroy the MCU on a rail fault that a current-limited supply would have caught for
nothing.

Bench supply on \file{VBAT} (\file{J2.2}), \textbf{16.8 V, current limit 100 mA}.

\begin{center}
\begin{tabular}{@{}cP{0.19\textwidth}P{0.22\textwidth}P{0.44\textwidth}@{}}
\toprule
\textbf{\#} & \textbf{Check} & \textbf{Expect} & \textbf{If wrong} \\
\midrule
1.1 & current at power-on & \textbf{$<$ 80 mA} & something is shorted: kill it and inspect
\file{U8}, \file{U18} and the bulk caps \\
1.2 & \file{+5V} rail & \textbf{4.97 V} $\pm$0.2 & \file{U8} TPS54202, \file{L1}, feedback divider \\
1.3 & \file{+3V3} rail & 3.3 V $\pm$0.1 & \file{U9} AP2112 \\
1.4 & \file{+3V3A} rail & 3.3 V $\pm$0.1 & \file{U10} TLV75533 \\
1.5 & \file{+9V} rail & \textbf{not populated} & \file{U18} is DNP by design \\
\bottomrule
\end{tabular}
\end{center}

Raise the limit to 500 mA only once every rail is correct.

\begin{quote}
\textbf{If a rail is dead with a correct supply, look at \file{R46} before the buck.} It is
\file{Q4}'s gate pull-down. Without it the gate floats, the FET never turns on, and the
board is dead with everything wired correctly.
\end{quote}

\begin{quote}
\textbf{Reverse polarity is protected, but do not test it.} \file{Q4} (WST4041 P-FET)
blocks a reversed pack and \file{DZ1} clamps the gate, but \file{D1} is still a TVS rather
than a series diode. Use the XT60 keying and check twice.
\end{quote}

\begin{quote}
\flag{USB cannot power this board.} VBUS runs the USB PHY only; the rails come from VBAT.
That would need a respin. For bench work, supply VBAT even to configure it.
\end{quote}

% =============================================================================
\section{Part 8 --- T2: MCU alive}

\begin{enumerate}
\item \flag{Bootloader over SWD first, before USB does anything.} Flash the board's own
bootloader from \file{Tools/bootloaders/} with \file{dfu-util}, using \file{TP20} (GND) and
\file{TP21} (+3V3) as the probe's reference. A board with no bootloader will not enumerate.
\item Data-capable USB-C cable --- charge-only leads enumerate nothing. Keep the supply
connected.
\item Hold \file{BOOT}, tap \file{RESET}. \file{lsusb} should show \file{0483:df11}.
\item Flash \file{arducopter.apj} from \file{tools/build\_firmware.sh}, pinned to
Copter-4.7.0.
\item Connect a ground station; confirm the board ID and read \textbf{every} boot message
--- every fitted sensor should probe, and nothing unfitted should be claimed.
\item \flag{Copy \file{copter-deadreckon-home.lua} to the microSD at \file{APM/scripts/} by
hand.} The build script does not package it into the \file{.apj}. Without it,
\file{FS\_EKF\_ACTION 0} ships with no GPS-loss response at all.
\end{enumerate}

% =============================================================================
\section{Part 9 --- T3: sensors on the bench, props OFF}

\begin{longtable}{@{}cP{0.21\textwidth}P{0.33\textwidth}P{0.33\textwidth}@{}}
\toprule
\textbf{\#} & \textbf{Check} & \textbf{How} & \textbf{Note} \\
\midrule
\endfirsthead
\toprule
\textbf{\#} & \textbf{Check} & \textbf{How} & \textbf{Note} \\
\midrule
\endhead
\bottomrule
\endfoot
3.1 & IMU1 (\file{U2}) orientation & roll/pitch, watch the horizon & \textbf{asserted in hwdef,
never measured} \\
3.2 & IMU2 (\file{U3}) orientation & same, \file{INS\_USE2} & only if \file{U3} was fitted \\
3.3 & Barometer & tracks a 1 m lift & \\
3.4 & Battery voltage & against a meter & \file{BATT\_VOLT\_MULT 11.000}, should be close \\
3.5 & \textbf{Battery current} & known load & ships \textbf{\file{BATT\_AMP\_PERVLT 25.0}} --- a
starting point, not a calibration \\
3.6 & RC input & \file{RC\_PROTOCOLS 512} (CRSF) & set to 1 for SBUS/FPort \\
3.7 & ESC telemetry & RPM with \file{SERIAL5\_PROTOCOL 16} & receive-only, and inert on a
BLHeli\_S ESC \\
3.8 & \textbf{Motor order and direction} & motor test, \textbf{props off} & \file{J2.3}--\file{J2.6}
= M1--M4 \\
3.9 & Buzzer, LEDs & arming tone, strip lights & \\
3.10 & TFS20-L, if fitted & tape measure at 1 / 5 / 10 m & \file{RNGFND1\_TYPE 46} \\
\end{longtable}

\textbf{Everything optional ships disabled in \file{defaults.parm} and must be turned on
here.} That is deliberate: a rangefinder configured but absent fails prearm, and a bare
board arms out of the box. Do not enable a sensor you have not fitted.

% =============================================================================
\section{Part 10 --- T3a: thermal, and the board measures half of it for you}

\file{U19} (TMP119) measures \file{U9} for you. It sits 3.9 mm from \file{U9} on I2C1 at
\file{0x48}, and \file{defaults.parm} ships \file{TEMP\_LOG 1}, so every flight writes a
\file{TEMP} message to the card.

\begin{center}
\begin{tabular}{@{}p{0.20\textwidth}p{0.42\textwidth}p{0.30\textwidth}@{}}
\toprule
\textbf{Part} & \textbf{Why it matters} & \textbf{Status} \\
\midrule
\file{U9} (3.3 V LDO) & \textbf{linear} regulator dropping 1.7 V in SOT-23-5 with no thermal
pad, carrying the MCU, flash and microSD & junction \textbf{104--157\degC} against a
\textbf{150\degC} limit. It has thermal shutdown --- being wrong means the 3.3 V rail
switching off in flight, not smoke \\
\file{U8} (+5 V buck) & \textbf{16.8 V $\rightarrow$ 5 V at 0.95 A} through a SOT-23-6 with
no thermal pad, sandwiched between a 60 A ESC and a battery & junction
\textbf{70--139\degC} against a \textbf{125\degC} \emph{recommended} max ---
\flag{an opinion until a thermocouple says so} \\
\bottomrule
\end{tabular}
\end{center}

Both brackets exist because the junction-to-ambient figure is a property of \emph{your}
assembled board and the two ends come from different sources (JEDEC and EVM columns of the
datasheet).

\textbf{Half of that bracket is now measured rather than assumed.} A SOT-23 has no thermal
pad, so every watt leaves through the leads into whatever copper is attached to them --- and
\file{tools/thermal\_vias.py} reads that copper off the board for \emph{both} regulators.
Measured 2026-09-19: \file{U8.1} has \textbf{3 vias within 1.7 mm holding 7655.8 mm$^2$ of
GND plane}, \file{U9.2} has 2 vias onto the same plane, and \textbf{no legal position
exists for another via on either part} (the search runs to 1.7 mm and refuses in-pad vias,
which would wick solder). That is the strongest statement the layout alone can make, and it
is why the low end of the bracket above is the plausible one for this board.
\flag{It still does not license a number} --- copper moves the bound monotonically, and only
a thermocouple converts it into degrees.

\textbf{What to do:} power it stack-assembled, logging to the card, ten minutes at steady
state. Read \file{TEMP[0]}. If you own a probe, put it on \file{U8} too and on \file{U9}'s
case once --- \file{U19} is 17.2 mm from \file{U8} and cannot speak for it, and one case
reading establishes the offset that \file{U19} then carries forward alone.

\textbf{Stop and ask if \file{TEMP[0]} exceeds about 85\degC{} or \file{U8}'s case exceeds
about 100\degC.}

% =============================================================================
\section{Part 11 --- T3b: the RF chain, and the single open question}

This is the one place the project is honest about not knowing.

\textbf{The Doppler solve is written and validated.} \file{tools/soop\_solver.py} runs
\file{SOLVER OK}, recovering a known position from clean geometry at \textbf{160.1 m median
at 5 Hz}, and reports a singular geometry rather than inventing a number.

\textbf{What is not known is what the shipping antenna and receiver actually deliver.} The
shipped \file{DopplerErrorModel} sigma is \textbf{180 m}; the measured divergence boundary
was taken at \textbf{20 m}. The SITL scenarios print this warning beside their own results:

\begin{quote}
\emph{``The 120 m divergence boundary DOES NOT APPLY here: it was measured at sigma 20 m
and this run is at sigma 180 m. Re-measure the boundary at the current sigma before reading
a hold/runaway verdict off it.''}
\end{quote}

\flag{Rescaling it by arithmetic would invent a measurement.} Only a bench RF measurement
moves this: a reference receiver, the real antenna, real geometry. Until then, do not read
a position-hold verdict off those scenarios.

% =============================================================================
\section{Part 12 --- T3c: IMU temperature calibration}

Free, and it attacks the dominant error term. Do it once the board is warm enough to be
realistic --- \textbf{4S, stack assembled}, not the board in free air, because the assembly
\emph{is} the thermal environment.

% =============================================================================
\section{Part 13 --- T3d: the nine unlabelled pads}
\label{sec:T3d}

Of the \textbf{30} pad and test-point references that exist on the board, \textbf{21} get a
silkscreen label. These \textbf{9} have no clean position at the board's own 0.80 mm
minimum text height, so \file{tools/silk\_labels.py} leaves them \textbf{blank rather than
printing something illegible} --- a wrong label on a board is worse than a blank pad,
because it will be trusted:

\begin{center}
\file{P41}\quad \file{P44}\quad \file{P62}\quad \file{P72}\quad \file{PL3}\quad
\file{TP4}\quad \file{TP6}\quad \file{TP7}\quad \file{TP9}
\end{center}

\flag{\file{P41} is 5 V.} \file{TP6} is \file{USART1\_RX}. Never identify a gold square by
counting from its neighbour --- use the printed gerber and \file{docs/img/padmap.svg}. The
gate reprints this list on every run, so treat any other number you read as stale.

% =============================================================================
\section{Part 14 --- T4/T5: tethered, then flight, strictly in this order}

\begin{enumerate}
\item \textbf{Run the SITL scenarios first:} \cmd{bash sitl/run\_scenarios.sh}. It found
three prearm faults no build check saw.
\item \textbf{Battery via the smoke stopper and the EC5-to-XT60 adapter} the first time.
\item \textbf{Tethered hover in Stabilize.} Land. Pull the log.
\item Tune \file{INS\_HNTCH\_FREQ} and \file{INS\_HNTCH\_BW} from the log's FFT.
\item \textbf{GPS Loiter.} This is a working aircraft.
\item \textbf{Props on last}, after motor direction is confirmed.
\end{enumerate}

\textbf{Everything past GPS Loiter is research.} GNSS-denied position hold is
uncharacterised until T3b is measured.

% =============================================================================
\section{Part 15 --- Audit: what is still manual, soldered, or unmeasured}

This is the audit you asked for. It is split into the three kinds of risk, because each is
removed by a different action: \emph{manual work} by buying a jig or an alternative,
\emph{hand soldering} by a stencil or an assembly tier, and \emph{unmeasured values} only by
a bench session.

\subsection{15.1 Manual soldering}

The board is machine-assembled, so the soldering is deliberately minimal. What remains:

\begin{center}
\begin{tabular}{@{}p{0.22\textwidth}cp{0.46\textwidth}@{}}
\toprule
\textbf{Item} & \textbf{Joints} & \textbf{Risk and mitigation} \\
\midrule
Buzzer, \file{PZ1}/\file{PZ2} & 2 & low. Must be a \textbf{passive} piezo; fitting an active
one costs you every arming and failsafe tone pattern, which you will only notice in the air \\
LED strip, \file{PL1}/\file{PL2}/\file{PL3} & 3 & low, and optional. \file{PL3} is unlabelled
--- see Section~\ref{sec:T3d} \\
\file{U3}, second IMU & 14-pin LGA & \flag{only if you chose Economic.} An LGA-14 has no
accessible leads, so it needs hot air or a stencil, not an iron. Standard assembly fits it
and removes this entirely \\
\bottomrule
\end{tabular}
\end{center}

\textbf{The inefficiency, stated plainly:} choosing Economic to save one part moves an
LGA-14 onto your bench and does \emph{not} save the feeder fees. It is the single manual
soldering task on this board that needs real equipment, and it buys you one optional IMU.
Unless you want the redundancy, Standard is the cheaper answer once your time is counted.

There is \textbf{no stencil requirement} and no hand-placed passives. That is by design.

\subsection{15.2 Manual assembly}

\begin{center}
\begin{tabular}{@{}P{0.25\textwidth}P{0.63\textwidth}@{}}
\toprule
\textbf{Task} & \textbf{Why it is manual, and what would remove it} \\
\midrule
The stack & four bolts, four spacers, grommets, the ESC and the frame's press nuts. Not
automatable and not worth automating \\
Motor screws & 16 screws whose length depends on a tapped depth you must measure. The M3
assortment removes the \emph{reorder}, not the measurement \\
Skids & printed, not bought. This is a feature: the 3.5 mm thickness is yours to set, which
is why the screw length is not a fixed number \\
Board connectors & the GPS, RC, I2C, lidar and ESC leads all plug in; no crimping except
the spare pigtails \\
Bootloader over SWD & four wires or pogo pins to a bare board. \flag{There is no
factory-installed bootloader, so this cannot be skipped} \\
The GPS-loss applet & \file{copter-deadreckon-home.lua} must be copied to the microSD by
hand; the build does not package it into the \file{.apj} \\
Companion mount & no mount exists yet. Deferred with the companion, so nothing blocks \\
\bottomrule
\end{tabular}
\end{center}

\textbf{Two structural constraints worth naming, because neither is removable now. Note
that the first reads like a defect and is not --- it is the opposite, and worth reading for
how it was settled:}

\begin{itemize}
\item \file{J3}'s GPS clearance \textbf{passes, because the requirement was re-derived
rather than the board changed.} The clear run before \file{R21} is \textbf{5.30 mm}. The
table's old figure of 6.0 mm was a JST-GH \emph{class} nominal, and it assumed a bend
allowance that nothing in that corridor needs: a pre-wired GH shell exits its wires at the
rear, already parallel to the board, and the only neighbour is a 0402 resistor 0.55 mm tall.
Shell half-width 2.6 mm plus 1.0 mm of handling margin is 3.6 mm, so 5.30 mm clears by
1.7 mm. \textbf{Relocating \file{R21} was mapped and rejected} --- \file{TP20}'s courtyard
and a via fence leave no position in the span that yields 6.0 mm, so there was never a fix
available there. The residual risk is a bulky moulded boot, which is why a spare JST-GH lead
is on the order list and T2 fits the plug \emph{before} the headers are soldered.
\item \flag{The board is 45.1 $\times$ 46.1 mm, below JLCPCB's 70 $\times$ 70 mm Standard
minimum.} Confirmed against their capability page: Standard single-PCB assembly runs
70 $\times$ 70 mm to 460 $\times$ 500 mm, while Economic runs 10 $\times$ 10 mm up. So
Standard needs the board panelled with breakaway rails --- an extra fabrication step and an
extra thing to confirm in their preview, not a blocker. Economic avoids the panel entirely
but charges the feeder fees anyway and drops an LGA onto your bench (Section 15.1), which is
why Standard plus panelisation is the recommendation.
\end{itemize}

\subsection{15.3 Measurements that have not been confirmed}

This is the longest list, and every row is a value some part of the build \emph{assumes}.
The gate's own verdict names nine of them as \file{NOT READY TO FLY}.

\begin{longtable}{@{}p{0.26\textwidth}p{0.34\textwidth}p{0.30\textwidth}@{}}
\toprule
\textbf{Unconfirmed value} & \textbf{What it decides} & \textbf{Confirmed by} \\
\midrule
\endfirsthead
\toprule
\textbf{Unconfirmed value} & \textbf{What it decides} & \textbf{Confirmed by} \\
\midrule
\endhead
\bottomrule
\endfoot
Motor boss tapped depth & M3 $\times$ 14 vs 12 + washer. Too deep a screw clamps nothing while
feeling tight & your calipers, before the first of 16 screws. Covered by the assortment \\
Press-nut tapped depth (30.5 mm joint) & M3 $\times$ 30 vs 28 + washer, on 6.7 mm of
engagement & your calipers \\
Skid thickness, \textbf{assumed} 3.5 mm & sets the motor screw length entirely & you, because
you print them. Measure the print \\
\textbf{Grommet flange} $\leq$ 5.9 mm & at 6.0 mm it touches a printed label; a 7 mm washer
overlaps by 0.52 mm & calipers on the grommets you receive \\
\textbf{ESC cable continuity} & whether \file{J2} matches the ESC you were sent. Wrong here
means a dead board or a smoked one & a multimeter, both ends, \textbf{before any VBAT} \\
\file{J3} plug against 5.30 mm & whether the received GPS lead seats & fit the plug before
the headers are soldered \\
\file{BATT\_AMP\_PERVLT} & battery failsafe and consumed-mAh. Ships 25.0; it is a property of
the ESC's shunt, not of this board & a known load at T3.5 \\
IMU orientations & \file{ROTATION\_ROLL\_180} is asserted in the hwdef, never measured. A
mirrored axis means it does not fly & roll and pitch it, watch the horizon (T3.1) \\
\file{U8} junction temperature & 70--139\degC{} modelled against a 125\degC{} recommended
max. If it really is at 139 the fix is a \textbf{respin}, not a reorder & a thermocouple, T3a \\
\file{U9}'s board-to-junction offset & makes \file{U19}'s logged value meaningful rather than
a nearby number & one probe reading on the case, then \file{U19} carries it \\
\file{FLOW\_ORIENT\_YAW} & the sign of the flow correction, if flow is ever fitted & bench
calibration \\
The RF sigma the antenna delivers & \textbf{everything} about GNSS-denied hold. Model says
180 m; the boundary was measured at 20 m & bench RF, T3b. \textbf{The single open question} \\
Whether this antenna can produce a usable Iridium fix at all & the project's premise & T3b \\
Flow quality over real grass & simulated flow is perfect flow & a field test \\
Radxa Zero 3W hole pattern & the companion tray, if the companion is ever fitted & \textbf{not
published} --- measure the board, then 3D-print a tray \\
Frame top-plate standoff spacing & nothing to buy: the kit supplies them & the kit \\
MicroSD withdrawal clearance & how you mount it & your mount \\
Grommet compression under load & a property of the rubber & the assembled stack \\
\bottomrule
\end{longtable}

\subsection{15.4 The three worst, ranked}

\begin{enumerate}
\item \flag{The RF sigma (T3b).} It is the only item that can invalidate the project's
premise rather than a build step, and no amount of desk work substitutes for it. Everything
downstream of it --- position-hold claims, the SITL scenarios' hold/runaway results --- is
explicitly not readable until it is measured.
\item \flag{\file{U8}'s junction temperature.} 139\degC{} against a 125\degC{} recommended
maximum is a measure-and-stop gate, and the remedy if it is real is a board respin. The
mitigation is that six layers with planes should land near the EVM end of the bracket ---
but that is an opinion until a probe says otherwise.
\item \flag{The ESC cable and the tapped depths.} Cheap to check, catastrophic to skip. Both
are on this list only because they are properties of parts somebody else manufactured, and
both are silenced by owning the M3 assortment and a multimeter.
\end{enumerate}

\subsection{15.5 Design-level inefficiencies, for a future revision}

These do not affect this order. They are recorded because a respin is the only place they
can be fixed, and one of them (\file{U8}) may force one anyway.

\begin{itemize}
\item \textbf{USB cannot power the board.} VBUS runs the PHY only. Every bench session needs
a VBAT supply, which is why T1 exists at all.
\item \textbf{Nine pads arrive unlabelled} (Section~\ref{sec:T3d}) because a three-character
label at the 0.80 mm minimum height covers a 2.84 mm pitch pad and its neighbours. More
silkscreen real estate, or fewer test pads, would remove a bench session.
\item \textbf{14 DNP designators still have pads and copper.} Correct for reversibility, but
it is unused copper on two rails.
\item \textbf{\file{-nofpv} is byte-identical to the default pair.} A variant that ships the
same bytes as another is a place for someone to pick the wrong file.
\item \textbf{The Economic variant saves one part and not the feeder fees}, which makes it a
worse deal than it reads as (Section 15.1).
\item \textbf{In4 rail fragmentation} is an accepted warning on this board.
\item \textbf{The 1620 MHz SAW filter is not stocked at LCSC}, so the RF front end needs a
second supplier and a separate board.
\end{itemize}

% =============================================================================
\section{Part 16 --- What cannot be confirmed from a desk}

Not everything here is improvable. These are genuinely outside what any local check or any
web lookup can settle, and they are listed so that ``verified'' is not read into them.

\begin{center}
\begin{tabular}{@{}p{0.40\textwidth}p{0.52\textwidth}@{}}
\toprule
\textbf{Item} & \textbf{Status} \\
\midrule
JLCPCB's upload preview & outline, panel size, CPL orientation, stackup, drills. No local tool
substitutes for looking at their render \\
1 oz outer copper & every current figure assumes it. Confirm what ships \\
LCSC stock & snapshot, age-gated at 7 days. Re-fetched live 2026-09-19: 56/56 codes resolve \\
The 1620 MHz SAW filter & not stocked at LCSC --- separate board, separate supplier \\
ArduPilot board ID & 7180 is free and below the ceiling; the request is written and ready.
9001 was free but above the ceiling the registry asks you to respect \\
\file{U8} junction temperature & modelled, not measured \\
The sigma the antenna delivers & \textbf{the single open question} \\
Real flow over grass & simulated flow is perfect flow \\
IMU orientations & asserted in the hwdef, never measured \\
\file{BATT\_AMP\_PERVLT} & a property of the ESC's shunt, not of this board \\
ESC cable continuity & buzz the cable you received \\
Radxa mounting holes & not published by the vendor \\
\file{J3} plug clearance & 5.30 mm clears the derived 3.6 mm requirement; what is still
unconfirmed is whether the plug you receive seats in it \\
Grommet flange & above 5.9 mm it touches a printed label \\
\bottomrule
\end{tabular}
\end{center}

% =============================================================================
\section{The one-page order}

\begin{quote}\small
\begin{enumerate}[label=\arabic*.,leftmargin=1.4em]
\item \cmd{python3 tools/preflight.py} --- must print \good{READY TO ORDER}
\item \cmd{bash tools/make\_order\_bundle.sh} --- regenerate everything
\item Upload \file{fab/NAVCORE-SoOP-order-bundle.zip}, or the gerbers plus the default
BOM/CPL pair
\item 6 layers, 45.1 $\times$ 46.1 mm, 1.6 mm, \textbf{1 oz} outer, ENIG, 4 mil, via
0.45/0.20
\item 5 bare boards, 2 assembled. Standard panelised 2 $\times$ 2 unless you want to
hand-fit \file{U3}
\item Confirm \file{Y1} = \file{C2682774} and \file{Y2} = \file{C5563878}; refresh the stock
snapshot if it is over 7 days old
\item Order the aircraft parts, \textbf{including the M3 assortment} and the Tier 4 RF kit.
Do not buy a PMW3901 or VL53L1X
\item While waiting: prepare the board ID request, print the padmap, decide the USB edge
\item On arrival: inspect under magnification, buzz the ESC cable, measure the tapped depths
\item T1 current-limited rails $\rightarrow$ T2 bootloader over SWD $\rightarrow$ T3 sensors,
props off
\item T3a thermal (\file{U8} needs a probe) $\rightarrow$ T3c IMU temp cal $\rightarrow$ T3d
nine unlabelled pads
\item T3b before any position-hold claim: measure the real sigma
\item SITL scenarios $\rightarrow$ tethered hover $\rightarrow$ GPS Loiter. \textbf{Props on
last}
\end{enumerate}
\end{quote}

\vfill
\begin{center}
\footnotesize\color{sidegrey}
Generated from the NAVCORE-SoOP working tree, 19 September 2026.\\
Gate: READY TO ORDER, 43 of 43 gated prerequisites proven.\\
Hardware: NOT READY TO FLY, 9 items requiring hardware.
\end{center}

\end{document}
